\section{Transmission: Factors and Lead--Lag Structure}
\label{sec:transmission}

% Arc: the transmission story. First the 20 factors, derived naturally; then
% the Cucuringu lead-lag features developed on top of them; ending with the
% final 4-block ridge of everything.
% Source material: alpha_manifestation_findings_2026-08-04.md,
% reference/sections/replication_state.tex (transmission width, span law).

The three blocks assembled so far are all \emph{contemporaneous}: the target's
own history, the exogenous panel, and the frozen pairwise products read the
current state of the market. This section adds the one channel that is not ---
directed lead--lag structure among the panel's latent factors, the
``transmission'' of activity from one part of the market to another at a
one-bar delay. The construction has three steps, and each falls out of the
data rather than out of a modelling choice: a factor frame that is nothing but
the panel's own first-window eigenstructure; a lead--lag field on that frame,
built and validated with the toolkit of Cucuringu and
coauthors~\pending{cite: Hermitian lead--lag clustering / SPONGE / dynamic-network
CPD references}; and finally the arrows entering the ridge as a fourth block.

\subsection{Twenty Natural Factors}

The factor construction involves no rotation, no supervision, and no tuning.
Take the panel's base features (of order 133 columns; 106 of them live,
non-degenerate directions), form their correlation matrix on the \emph{first
training window} (2005--06), and project the panel onto its leading
eigenvectors. The frame is then \textbf{frozen permanently}: factor scores
$G_i(t)$ for the rest of the sample are projections onto directions estimated
once, from feature data alone, before any forecasting begins. Because the
frame is estimated from $X$ and never from the target, the geometric objects
built on it can be validated split-half against the signal without
contamination --- which is precisely why they turn out to be the study's most
replicable objects.

In what sense is the count natural rather than chosen? Three independent
statistics of the same first-window spectrum bracket it. The participation
ratio of the spectrum is $\approx 18$ --- an effectively $\approx 20$-dimensional
span. The raw Marchenko--Pastur edge (iid null at the nominal sample size)
keeps 22 of the 106 live eigenvalues above noise. And the frame itself is the
one linear object in the study that is stable across halves of an 18-year
sample --- the factors' \emph{nonlinear} self-geometry (curvature, spectra of
interaction forms, cluster partitions) fails every split-half test, while the
linear frame holds. Twenty is the round cover of an 18--22 bracket that the
data produces on its own; it was carried as a convention, not optimized, and
Section~\ref{sec:trans-leadlag} reports the arm that later stress-tested it.

One refinement of the Marchenko--Pastur picture matters for everything that
follows. Bar data is heavily dependent: the factor scores' median integrated
autocorrelation time is 175 bars ($\approx 3.6$ days), which deflates the
effective sample to $T_{\mathrm{eff}} \approx 68$ per dimension and raises the
dependence-adjusted noise edge to $5.04$ --- above which only \textbf{5}
eigenvalues survive as individually identified directions. Read correctly,
this is not ``the mid-spectrum is noise'' (objects built on it replicate
split-half at $+0.6$ to $+0.99$, which noise cannot do); it is that
mid-spectrum eigenvectors are \emph{gauge-degenerate}: the data pins the
subspace they span but not the basis inside it. The factors are a stable
\emph{span} with unpinnable labels. This one fact predicts the section's
central robustness result --- that any generous cover of the span works
equally well, and no clever selection inside it can beat carrying it whole.

\subsection{Cucuringu Lead--Lag Features}
\label{sec:trans-leadlag}

\paragraph{The lead--lag field is real, and it is a rotation.}
Let $A$ be the antisymmetric part of the factor scores' lag-one-day
cross-correlation --- the raw material of Hermitian lead--lag
clustering~\pending{cite}. Over its 190 edges, $A$ replicates split-half at
$+0.79$: the second geometric object on the frame (after the factor-pair
coupling map, split-half $+0.62$ against a null of $0.00 \pm 0.11$) to pass a
powered replication test. Its structure is diagnosed by the Hodge
decomposition: \textbf{91.6\% curl, 8.4\% gradient}. The daily field is not a
lead--lag \emph{hierarchy} but a \emph{circulation} --- one dominant two-plane
carries 61\% of the antisymmetric energy, and the Hermitian phase ordering of
the factors around that cycle replicates across halves at $+0.76$. The
toolkit adjudicated its own instruments symmetrically: methods that presume a
ranking (SyncRank: split-half $\tau = -0.06$) or a partition (SPONGE signed
clustering: split-half ARI below its permutation null at every $k$) fail on
this object \emph{because} it is curl-dominated and dense-but-weak has no
blocks, while the Hermitian embedding and dynamic-network change-point
detection pass --- the latter flagging exactly two interior breaks (2013-07,
the taper transition, and 2017-08, immediately preceding Volmageddon) with no
tuning toward them, and, notably, \emph{no} break at COVID: amplitudes
exploded in 2020 while the field's structure held.

\paragraph{At the product's own timescale.}
The same construction at a one-\emph{bar} lag --- with per-slot
deseasonalization and an explicit clock control, since the diurnal timetable
could manufacture artificial lead--lag --- yields the most replicable object
in the study: split-half $+0.992$ for the dynamic flow (the clock share is
$-0.003$; a closed-form control for trailing-MA arithmetic explains 6\% of
its variance, and the residual still replicates at $+0.992$). The intraday
field is even more cyclic (96.5\% curl), operates identically in and out of
regular hours (RTH-vs-overnight similarity $+0.967$), and its hub structure
is economically legible: turnover and tail activity lead the implied-vol
complex, and implied vol leads the realized-vol factors --- order flow
repricing implied variance, implied variance leading realized, at a
half-hour lag.

\paragraph{From field to feature.}
The feature construction is the field applied as a linear operator. With
$D$ the antisymmetric part of the lag-one-bar cross-correlation, estimated
causally on a trailing 504-day window with quarterly re-estimation, the
transmission columns are
\[
\hat G_j(t) \;=\; \sum_i D_{ij}\, G_i(t - 1\ \mathrm{bar}),
\]
one column per factor: each direction's next-bar motion as predicted by the
flow through it. The antisymmetry is load-bearing, not decoration --- the
symmetric part of the lag-1 cross-correlation is persistence content the
moving-average ladder already spans. That claim was given its decisive
control: 20 columns of \emph{plain} lagged scores $G(t-1)$ (the same lagged
information with no flow structure) fail the entry gate outright ($+0.43$),
while the $D$-weighted columns pass ($+2.58$ at the eight-bar horizon). The
flow content, not generic lag-1 information, is what the design was missing.

The block's value is fast and strictly horizon-dependent. Against same-engine
twins that already contain everything else, the Diebold--Mariano profile is
$+9.90$ at one bar (the strongest feature statistic in the campaign),
$+5.00$ at four bars, $+2.58$ at eight, failing at the end-of-day horizon
($+1.38$), with one lawful revival at sixteen bars ($+5.85$) --- but only
under a per-horizon shrinkage scaling $\alpha \propto H$ that prices the
block at sixteen times its one-bar penalty. Dense-weak content pays at long
horizons only when the solver is forced to be humble. One further symmetry:
the arrows themselves do not want freshness. Re-estimating $D$ monthly
($+0.75$) or even per bar ($+1.48$) fails the gate --- the refit-cadence law
that governs coefficients everywhere else in this paper finally met the
object it does not govern, and it is precisely the study's most stable
object. Lag-one $D$, re-estimated quarterly, is the complete form of the
feature.

\paragraph{The width plateau.}
``How do we know it is just 20?'' --- the honest answer on record was
\emph{we don't}: the width was inherited. Varying it end to end, with the
frames nested (each width $q$ takes the top-$q$ eigenvectors of the same
first-window correlation), produced the campaign's cleanest structural
result. The one-bar increment is not flat in $q$: it improves through
$q = 40$ (head-to-head over $q=20$: $+2.29$ at daily refit, $+4.59$ per-bar),
holds a shoulder through $q = 80$, and fails its gate outright at $q = 106$
--- every live direction ($+1.62$ against the twin, the width ladder's first
failure). Refinement rungs at $q \in \{30, 50, 60, 70\}$ show the interior is
flat to within noise. The forecast improvement sits on a \textbf{plateau over
widths $\sim$30--80 with cliffs on both sides}: the arrows live in a
finite-rank sub-span roughly \emph{twice} the conventional frame, dense
within it and empty beyond it. Production runs width 40 as a representative
of the plateau, not as a tuned optimum --- exactly the robust-region shape
the gauge-degeneracy story predicts, since covering the span is what matters
and its truncation point is not identified. The mechanics were probed
directly: the widening gain is carried by the new directions' \emph{own}
next-bar forecastability (not by new information about the old factors);
doubling the $D$ estimation window rescues nothing at $q=106$ (the tail is
empty, not noisy); the new arrows replicate split-half at $+0.838$ (old
block $+0.907$) --- a stable field over unpinnable labels; and the width
increment is horizon-local, existing at the one-bar horizon only (the
width-40-over-width-20 gate fails at $H = 4$, $8$, $13$, and $16$), which
retroactively explains why the inherited 20 survived so long: at every
horizon tested earlier, 20 was enough.

\subsection{The Four-Block Ridge}
% Endpoint: the final composed model — all four blocks.

The transmission block completes the model. The final specification is one
blockwise ridge in a single solver --- \textbf{719 columns} --- refit every
bar on a 250-day rolling window with a one-day embargo, with the penalty
graded by block:
\[
y \;\sim\; \mathrm{ridge}_{250\mathrm{d}}
\bigl[\;
\underbrace{\text{target-HAR backbone}}_{27\ \text{cols},\ \alpha_{\mathrm{eff}} = 1}
\;\big|\;
\underbrace{\text{linear exogenous block}}_{516 + 36\ \text{cols},\ 3\times10^{3}}
\;\big|\;
\underbrace{\text{frozen products}}_{100\ \text{cols},\ 3\times10^{4}}
\;\big|\;
\underbrace{\text{transmission}}_{40\ \text{cols},\ 3\times10^{3}}
\;\bigr].
\]
\decide{penalty presentation: the campaign record states the blocks at
effective $\alpha = 1$ / $3\times10^{3}$ / $3\times10^{4}$ by column scaling,
while the arc notes for Sections~\ref{sec:dense_weak}
and~\ref{sec:nonlinear_product} quote the $1/100/1000$ normalization ---
harmonize one convention across the three endpoint specs}
The transmission block's own shrinkage was laddered and stays at the
exogenous penalty --- it behaves like exogenous content, not like products.
The backbone is unshrunk, the linear panel is shrunk hard, the quadratic
block an order of magnitude harder, and the arrows ride at the linear
penalty: collect everything, refresh constantly, grade the humility by
block.

The increment is small and certain, in the manner of every gain in this
paper. On the campaign panel, adding the width-40 transmission block to the
699-column per-bar model moves QLIKE from $0.12546$ to $0.12529$
(Diebold--Mariano $+4.59$; $+2.32$ on 2020+), against that panel's HAR
baseline of $0.13275$. Composed with the second-moment correction of
Section~\ref{sec:smearing} in its final form, the campaign's headline number
closed at QLIKE $0.12299$ ($\approx -7.3\%$ against HAR), having opened the
campaign at $0.12579$. Every step of the transmission program --- the one-bar
per-bar entry ($t = +8.45$), the width move, the intermediate-horizon rungs
--- survived a closing Romano--Wolf step-down over the full post-hoc family
(2{,}000 draws, 1{,}008-bar joint blocks), the study's first family sweep
with no casualties. \decide{which number is this section's endpoint headline:
the four-block ridge alone ($0.12529$ per-bar) or the ridge composed with
the Section~\ref{sec:smearing} smear ($0.12299$) --- and whether the paper's
final table re-freezes these on its own battery}
\pending{final four-block QLIKE regenerated under this paper's frozen battery
and panel convention, if the campaign numbers above are not adopted as-is}

The per-horizon deployment falls out of the horizon profile rather than out
of a choice: at one bar, the full four-block model refit per bar; at
intermediate horizons, the width-20 block with $\alpha \propto H$ scaling; at
end-of-day, no transmission at all --- the three-block model, by measurement
rather than omission. The postscript writes itself: the fourth block is the
first in the model with an arrow of time in it, and on this panel it is
descriptive at a one-day lag and predictive at a one-bar lag --- fast, like
everything else that worked. On a cross-section of assets, where the
lead--lag field is estimated across names rather than across one market's
factors, the same construction becomes the estimation target rather than the
garnish; that is the natural continuation of this section and lies beyond
this paper.
